\documentclass[10pt]{article}
\providecommand{\FIGDIR}{fig}
\newcommand{\DIGESTDATE}{14 September 2026}
% ===================== PM Research Watch - shared preamble =====================
\usepackage[T1]{fontenc}
\usepackage[utf8]{inputenc}
\usepackage{amsmath,amssymb}
\usepackage{tgpagella}
\usepackage{tgheros}
\renewcommand{\sfdefault}{qhv}
\renewcommand{\ttdefault}{lmtt}
\usepackage{microtype}
\usepackage[english]{babel}

\usepackage[a4paper,top=17mm,bottom=17mm,left=17mm,right=17mm,
            headsep=6mm,footskip=10mm]{geometry}
\usepackage{graphicx}
\usepackage{xcolor}
\usepackage{colortbl}
\usepackage{tikz}
\usetikzlibrary{calc}
\usepackage[most]{tcolorbox}
\usepackage{titlesec}
\usepackage{fancyhdr}
\usepackage{booktabs}
\usepackage{array}
\usepackage{longtable}
\usepackage{enumitem}
\usepackage{multicol}
\usepackage{ragged2e}
\usepackage[hidelinks]{hyperref}

% ------------------------------------------------------------------ palette
\definecolor{Ink}    {HTML}{14262E}
\definecolor{Deep}   {HTML}{0F4C5C}
\definecolor{Teal}   {HTML}{1F7A8C}
\definecolor{Sky}    {HTML}{4FA3C4}
\definecolor{Sage}   {HTML}{6FA287}
\definecolor{Amber}  {HTML}{E3A008}
\definecolor{Clay}   {HTML}{B5643C}
\definecolor{Coral}  {HTML}{D1495B}
\definecolor{Violet} {HTML}{6A5B9E}
\definecolor{Slate}  {HTML}{4A6572}
\definecolor{Paper}  {HTML}{FBF9F5}
\definecolor{Mist}   {HTML}{EEEAE2}
\definecolor{Rule}   {HTML}{D9D3C9}

\pagecolor{Paper}
\color{Ink}

\hypersetup{colorlinks=true, linkcolor=Deep, urlcolor=Teal, citecolor=Teal,
            pdftitle={PM Research Watch}, pdfauthor={Automated daily digest}}

% ------------------------------------------------------------------ headings
\titleformat{\section}
  {\sffamily\bfseries\large\color{Deep}}{}{0pt}
  {}[{\vspace{1.2mm}\color{Amber}\titlerule[1.1pt]}]
\titlespacing*{\section}{0pt}{5.5mm}{2.8mm}

\titleformat{\subsection}
  {\sffamily\bfseries\normalsize\color{Teal}}{}{0pt}{}
\titlespacing*{\subsection}{0pt}{3.6mm}{1.4mm}

\setlist[itemize]{leftmargin=4.2mm, itemsep=1.1mm, topsep=1.3mm, parsep=0pt,
                  label=\textcolor{Teal}{\footnotesize$\blacktriangleright$}}

% ------------------------------------------------------------------ header/footer
\pagestyle{fancy}
\fancyhf{}
\renewcommand{\headrulewidth}{0pt}
\renewcommand{\footrulewidth}{0pt}
\fancyhead[L]{\sffamily\scriptsize\color{Slate}PM RESEARCH WATCH\;\textcolor{Rule}{\textbar}\;\DIGESTDATE}
\fancyhead[R]{\sffamily\scriptsize\color{Slate}Particulate matter monitoring \& health effects}
\fancyfoot[C]{\sffamily\scriptsize\color{Slate}\thepage}

% ------------------------------------------------------------------ components
% metric tile
\newcommand{\metric}[3]{%
  \begin{tcolorbox}[enhanced, colback=white, colframe=Rule, boxrule=0.5pt,
    arc=1.6mm, left=2.2mm, right=2.2mm, top=1.8mm, bottom=1.8mm,
    borderline west={1.7mm}{0pt}{#3}, nobeforeafter, valign=center,
    height=17mm, width=\linewidth]
    \RaggedRight
    {\sffamily\bfseries\fontsize{17}{17}\selectfont\color{#3}#1}\\[0.5mm]
    {\sffamily\fontsize{6.5}{7.4}\selectfont\color{Slate}#2}
  \end{tcolorbox}}

% pull-quote / key insight box
\newtcolorbox{keybox}[1][]{enhanced, colback=Mist, colframe=Mist,
  arc=1.6mm, left=3.2mm, right=3.2mm, top=2.4mm, bottom=2.4mm,
  borderline west={1.7mm}{0pt}{Amber}, #1}

% subtopic band
\newcommand{\band}[2]{%
  \vspace{3.2mm}\noindent
  \begin{tcolorbox}[enhanced, colback=#2, colframe=#2, arc=1.2mm,
    left=2.6mm, right=2.6mm, top=1.1mm, bottom=1.1mm, boxrule=0pt,
    width=\linewidth, nobeforeafter]
    {\sffamily\bfseries\small\color{white}#1}
  \end{tcolorbox}\vspace{1.6mm}}

% one paper entry: {short cite}{title}{journal}{doi}{body}
\newcommand{\paperentry}[5]{%
  \par\vspace{1.9mm}\noindent
  {\sffamily\bfseries\footnotesize\color{Ink}#1}\;%
  {\sffamily\fontsize{7}{8}\selectfont\color{Slate}#3}\\[0.4mm]
  {\footnotesize\itshape\color{Deep}#2}\\[0.6mm]
  {\footnotesize #4}\\[0.3mm]
  {\fontsize{6.6}{7.6}\selectfont\color{Slate}\texttt{doi:}\,\href{https://doi.org/#5}{\texttt{\detokenize{#5}}}}
  \par}

% caption style
\newcommand{\figcap}[1]{\par\vspace{1.0mm}
  {\fontsize{7.2}{8.6}\selectfont\color{Slate}\textsf{\textbf{Fig.}}\; #1\par}\vspace{1.5mm}}

\providecommand{\permil}{\textperthousand}

\setlength{\parindent}{0pt}
\setlength{\parskip}{0pt}
\renewcommand{\arraystretch}{1.18}


\begin{document}

% ==================================================================== MASTHEAD
\thispagestyle{empty}
\begin{tikzpicture}[remember picture, overlay]
  \fill[Deep] (current page.north west) rectangle
      ($(current page.north east)+(0,-67mm)$);
  \fill[Amber] ($(current page.north west)+(0,-67mm)$) rectangle
      ($(current page.north east)+(0,-68.4mm)$);
  \foreach \x/\y/\r in {22/12/0.9, 41/26/1.5, 63/9/0.7, 88/31/1.1, 112/17/1.9,
                        138/29/0.8, 159/13/1.3, 178/54/1.0, 196/21/1.6,
                        30/44/0.6, 74/57/1.2, 125/48/0.9, 168/6/0.8, 191/59/0.7,
                        52/18/1.0, 100/60/0.6, 148/20/1.1, 66/49/1.4, 205/45/0.8}
    \fill[Sky, opacity=0.30] ($(current page.north west)+(\x mm,-\y mm)$) circle (\r mm);
\end{tikzpicture}

\vspace*{4mm}
{\sffamily\fontsize{7.6}{9}\selectfont\color{Sky}%
 AUTOMATED DAILY LITERATURE DIGEST\;\textbullet\;PhD RESEARCH WATCH\par}
\vspace{2.2mm}
{\sffamily\bfseries\fontsize{27}{29}\selectfont\color{white}%
 Particulate Matter\par}
{\sffamily\fontsize{27}{29}\selectfont\color{Sky}%
 Monitoring \& Health Effects\par}
\vspace{5.0mm}
{\sffamily\fontsize{8.4}{10}\selectfont\color{white}%
 Issue 2026-09-14\;\;\textcolor{Amber}{\textbar}\;\;PubMed entry window 14 September 2026\;\;%
 \textcolor{Amber}{\textbar}\;\;16 records screened in scope\par}
\vspace{18mm}

% -------------------------------------------------------------- metric strip
\noindent
\begin{minipage}[t]{0.190\linewidth}\metric{16}{RECORDS IN SCOPE\\(11 EXCLUDED, 6 HELD)}{Deep}\end{minipage}\hfill
\begin{minipage}[t]{0.190\linewidth}\metric{6}{OF 10 SUBTOPIC\\BANDS POPULATED}{Teal}\end{minipage}\hfill
\begin{minipage}[t]{0.190\linewidth}\metric{5}{SENSING RECORDS\\(31\% OF CORPUS)}{Sky}\end{minipage}\hfill
\begin{minipage}[t]{0.190\linewidth}\metric{4}{EFFECT ESTIMATES\\(3 WITH A CI, 0 ON PM)}{Amber}\end{minipage}\hfill
\begin{minipage}[t]{0.190\linewidth}\metric{8}{RECORDS CARRYING\\NO PMID}{Clay}\end{minipage}

\vspace{6mm}

% -------------------------------------------------------------- signal of day
\section*{\sffamily\bfseries\color{Deep}Signal of the day}
\vspace{-1.5mm}{\color{Amber}\rule{\linewidth}{1.1pt}}\vspace{2.5mm}

\begin{keybox}
{\sffamily\bfseries\small\color{Deep}Withdrawing a monitor changed the pollution, not just
the record of it.}\\[1.4mm]
\footnotesize
Zhou et al.\ (\textit{PNAS}) treat the 2025 suspension of the US embassy air-quality
programme as a natural experiment and ask what happens to a city when a trusted,
independent monitoring node goes dark. Using real-time satellite AOD in a matched
difference-in-differences design against same-country control cities, they find
\textbf{no differential pre-trend} across 46 cities and then a prompt
\textbf{24--33\% rise in local pollution} once reporting stops, largest and most
persistent where no substitute data source existed. Calibrated against concurrent
pre-shutdown embassy PM$_{2.5}$, the preferred conservative translation is
\textbf{+1.9\,$\mu$g/m$^3$}; alternative translations span 1.5 to
16.3\,$\mu$g/m$^3$, and no confidence interval is reported for the headline effect ---
so the causal design is considerably stronger than the magnitude. The result inverts
the usual framing of this watch. Dense low-cost networks are normally justified as a
way to \emph{observe} exposure at a resolution regulatory networks cannot reach; here,
observation is itself the intervention, and its withdrawal is measurable in the
atmosphere. Whatever a distributed network costs, this paper is the argument that its
governance and continuity matter as much as its calibration.
\end{keybox}

\vspace{3mm}

% -------------------------------------------------------------- exec brief
\section*{\sffamily\bfseries\color{Deep}What changed today}
\vspace{-1.5mm}{\color{Amber}\rule{\linewidth}{1.1pt}}\vspace{2.5mm}

\footnotesize
\begin{itemize}
\item \textbf{A thin day on the health axis, and the thinness is topical.} Registry-wide
PubMed deposits for the 14 September entry date stand at \textbf{2{,}375}, against 4{,}593
for 13 September. Both PubMed queries together returned eight candidates; two were
screened out and none reported a new PM exposure--response estimate. Four of the ten
subtopic bands --- \textit{Neuro}, \textit{Cardiovascular \& metabolic},
\textit{Reproductive \& developmental} and \textit{Other clinical} --- are empty, the
eighth issue of fifty to populate exactly six.

\item \textbf{Instrumentation carried the issue, and PubMed indexed almost none of it.}
Five of sixteen records are sensing or instrumentation work; \textbf{eight of sixteen
records carry no PMID at all}, arriving through the Crossref-by-ISSN leg
(\textit{AMT}, \textit{ACP}, \textit{Atmosphere}, \textit{EST}) or the Consensus sweep.
A PubMed-only watch would have shipped a six-record issue today.

\item \textbf{Regional, not universal, is now the settled answer on low-cost correction.}
Gupta et al.\ conclude from PA-II co-locations in Raleigh and Delhi that
location-specific pre-deployment calibration is a prerequisite for scientific use, and
G\"okozan's ablation study shows that where continuous co-location is not feasible,
humidity-interaction and cyclic-time features recover most of the gap
(test $R^2$ 0.724 $\to$ 0.847) but leave an RMSE of 5.19\,$\mu$g/m$^3$ --- large against
typical ambient means. The two papers agree on the diagnosis and differ only on the
fallback.

\item \textbf{No PM exposure estimate with a confidence interval was published into this
window.} The three interval-bearing estimates in today's forest plot are all
\emph{heat} exposure, from the COPD meta-analysis; the fourth, the embassy effect, has
no interval. This is the thirteenth issue in fifty with no PM ratio estimate carrying a
CI, so the figure should be read as a record of what was available, not as evidence
about PM.

\item \textbf{Two records were removed by the DOI gate after title screening passed
them.} Both Consensus hits on PurpleAir transferability and Bogot\'a LCS calibration had
already shipped --- on 29 July and 3 August --- under \emph{rewritten} titles, one a
paraphrase and one entirely different. Title-normalised deduplication is structurally
incapable of catching this; only the DOI check is.

\item \textbf{Six in-scope records were verified and deliberately not shipped.} Their
DOIs resolve, but no abstract exists in Crossref, Europe PMC, Scholar Gateway or
Consensus, and the publisher pages are client-rendered. The priority pickup is the
\textit{Aerosol Science and Technology} Spider-MAGIC and OPS size-distribution
evaluation --- a tier-A instrument paper held only for want of 200 words.
\end{itemize}

\clearpage

% ==================================================================== ANALYTICS
\section*{\sffamily\bfseries\color{Deep}Corpus analytics}
\vspace{-1.5mm}{\color{Amber}\rule{\linewidth}{1.1pt}}\vspace{3mm}

\noindent\includegraphics[width=\linewidth]{\FIGDIR/f1_subtopics.png}
\figcap{Subtopic assignment of the 16 in-scope records, single-label by dominant
contribution. \textit{Sensing, forecasting \& instrumentation} leads at 5, followed by
\textit{Exposure assessment \& modelling} at 4 and \textit{Occupational \& indoor} at 3.
Only six of the ten bands are populated; the four clinical-epidemiology bands that
usually carry a daily issue are all empty.}

\vspace{2mm}
\noindent\includegraphics[width=\linewidth]{\FIGDIR/f2_design_endpoint.png}
\figcap{Left --- study architecture. Modelling and inventory work is the largest class (5),
with measurement campaigns and reviews at 3 each and sensor co-location at 2; a single
record is observational cross-sectional and a single one chamber/laboratory. Sixteen
design labels were coined for this issue and added to the classifier in the same edit,
because the existing map was built for clinical epidemiology and this corpus is not.
Right --- endpoint. \textbf{Thirteen of sixteen records carry no human health endpoint}:
seven report air-quality concentrations, three emission flux, and three sensor or
instrument performance. Three clinical endpoints appear in total, one each ---
respiratory, cerebrovascular and all-cause mortality.}

\vspace{2mm}
\noindent\includegraphics[width=\linewidth]{\FIGDIR/f3_metric_geography.png}
\figcap{Left --- particle metric. Unspeciated PM or emissions is the exposure in 10 of 16
records, PM$_{2.5}$ alone in 3; composition/speciation and optical properties appear once
each. Right --- geographic provenance. Nine records chart as global or multi-region, which
reflects the day's composition rather than breadth: burden syntheses, satellite products
and multi-continent syntheses land there, and the purely numerical sMCPC study folds into
the same bucket for want of a geography. The Raleigh-and-Delhi co-location spans two
regions and is charted under one; China contributes a single record.}

\vspace{2mm}
\noindent\includegraphics[width=\linewidth]{\FIGDIR/f6_lifecourse.png}
\figcap{Life-course windows addressed. \textbf{Three of the five stages are empty today} ---
no preconception or in-utero, infancy or adolescent record appears at all. The
working-age count reflects two occupational monitoring deployments (three fire stations;
hospital staff areas), not a health-outcome cohort, and the 65+ entry is the COPD
meta-analysis, which does not report an age-stratified estimate.}

\vspace{2mm}
\noindent\includegraphics[width=\linewidth]{\FIGDIR/f4_heatmap.png}
\figcap{Subtopic $\times$ study-architecture cross-tabulation. The grid is sparse by
construction at $n=16$, and the structure it shows is the day's defining feature:
the sensor co-location cells sit entirely inside \textit{Sensing}, the modelling cells
spread across \textit{Sensing} (2), \textit{Exposure assessment} (2) and
\textit{Burden} (1), and no cell links an instrumentation architecture to a clinical
band --- the two clinical rows each hold a single review. Instrumentation and epidemiology did not
meet anywhere in this corpus.}

\vspace{2mm}
\noindent\includegraphics[width=\linewidth]{\FIGDIR/f5_forest.png}
\figcap{All quantitative effect estimates in the corpus on a common logarithmic axis ---
four, which is few, and \textbf{none of them is a particulate exposure}. Three are
\emph{heat} exposure from the 36-study COPD meta-analysis, shown here because the
air-pollution-adjusted sensitivity analysis is what makes them relevant to this watch;
the mortality and morbidity estimates are strata of the overall estimate and are not
independent of it. The fourth is the lower bound of the 24--33\% embassy-shutdown range,
plotted without an interval because none was reported. Read this figure as a statement
about what the window contained, not about PM risk.}

\clearpage

% ==================================================================== DIGEST
\section*{\sffamily\bfseries\color{Deep}Paper-level digest}
\vspace{-1.5mm}{\color{Amber}\rule{\linewidth}{1.1pt}}\vspace{1mm}

{\fontsize{7.4}{9}\selectfont\color{Slate}Ordered by cluster. Each entry gives the
methodological core and the finding that matters, not an abstract paraphrase.\par}

\band{Sensing, forecasting \& instrumentation\hfill 5 records}{Amber}

\paperentry{Zhou et al.\ 2026}
{Air pollution after the shutdown of US embassy monitoring}
{Proc Natl Acad Sci U S A}
{Matched difference-in-differences across 46 cities with same-country controls, using
real-time satellite AOD. No differential pre-trend; a prompt \textbf{24--33\% rise in
local pollution} after US reporting stopped, largest where substitutes were poorest.
Calibrating AOD to concurrent pre-shutdown embassy PM$_{2.5}$ gives \textbf{+1.9\,$\mu$g/m$^3$}
on the authors' conservative translation, with alternatives spanning 1.5--16.3 ---
an order-of-magnitude spread that sits in the calibration step, not the DiD, and no CI
accompanies the headline. The design is the contribution: monitoring treated as an
intervention with a measurable atmospheric effect.}
{10.1073/pnas.2605467123}

\paperentry{Zhou et al.\ 2026}
{Performance modeling of a small mixing-type condensation particle counter}
{Atmos Meas Tech}
{Validated CFD model of the sMCPC particle channel resolving turbulent flow, temperature
and species transport, with Kelvin, non-continuum and latent-heat terms in the growth
calculation. Raising the saturator--cooler temperature difference $\Delta T$ raises
supersaturation and lowers the critical activation diameter $D_{p,50}$; raising total flow
$Q_g$ suppresses supersaturation and shifts its peak downstream; both $\Delta T$ and
vapour fraction steepen the $D_{p,50}$ slope. Everything is numerical --- counting
efficiency is asserted through one validation case, not measured across the swept space.
This is design-space guidance of the kind a portable ultrafine network node needs
before it goes to the field.}
{10.5194/amt-19-5815-2026}

\paperentry{Gupta et al.\ 2026}
{Evaluating and calibrating low-cost air quality sensors in contrasting aerosol regimes}
{Aerosol Air Qual Res}
{PurpleAir PA-II co-located with FEM monitors for 2--4 weeks in Raleigh, North Carolina
and Delhi, India, deliberately spanning a clean and a heavily loaded regime. A model on
PA PM$_{2.5}$, temperature and RH holds mean absolute bias within \textbf{12\% hourly and
6\% daily}. The headline is negative: pre-deployment calibration specific to location or
region is a prerequisite for scientific use, i.e.\ no single transferable correction
exists. Two to four weeks is short against the seasonal cycle in either city, so drift
and seasonal composition change lie outside what this design can see.
\textit{Connector-sourced; outside the 14 September entry window.}}
{10.1007/s44408-026-00146-x}

\paperentry{G\"okozan 2026}
{Machine-learning calibration of a low-cost PM$_{2.5}$ sensor with location-specific
covariates and feature engineering}
{Chemosensors}
{Ablation design separating environmental covariates, auxiliary gases and feature
engineering for one outdoor residential sensor, validated on a leakage-safe time-based
split. Feature engineering helped in every configuration: raw PM$_{2.5}$ plus engineered
features raised Random Forest test $R^2$ from \textbf{0.724 to 0.847} and cut RMSE from
6.96 to \textbf{5.19\,$\mu$g/m$^3$}; humidity interactions and cyclic time encodings carried
the importance. One sensor at one site is the limit, and 5.19\,$\mu$g/m$^3$ is large
against typical ambient means --- but the operational claim, that useful correction is
reachable without continuous co-location, is the one network operators need tested.
\textit{Connector-sourced; outside the entry window.}}
{10.3390/chemosensors14070154}

\paperentry{Sun et al.\ 2026}
{Aerosol properties during a transatlantic smoke transport event from Aeolus observations}
{Atmos Chem Phys}
{ALADIN, a spaceborne high-spectral-resolution lidar, retrieves extinction and backscatter
directly rather than through an assumed lidar ratio, and is used to follow a
September 2020 western-US smoke plume to Europe. Mean extinction near 200\,Mm$^{-1}$
across selected cross-sections; AOD fell 0.54 $\to$ 0.33 along the path; the layer rose
from $\sim$4.4\,km over the US to above 6\,km over Europe and arrived above cloud, with
lidar ratios first falling then rising. One event constrains generality. For
satellite-to-surface PM$_{2.5}$ fusion the relevant point is blunt: this plume was
invisible at ground level throughout.}
{10.5194/acp-26-12953-2026}

\band{Exposure assessment \& modelling\hfill 4 records}{Amber}

\paperentry{Ajibade et al.\ 2026}
{Gentrification, displacement, disadvantage and exclusivity predict urban environmental
and health inequities}
{Proc Natl Acad Sci U S A}
{PCA on 2010 and 2020 ACS and Census data yields four contemporary segregation typologies
for 342 Portland tracts, compared against historical HOLC zones across 104 environmental
and health variables. HOLC grade predicted persistent exclusivity in the urban core but
not injustice elsewhere; \textbf{persistently disadvantaged tracts, often outside HOLC
boundaries entirely}, carried the highest pollution and hazard exposure and the worst
health and disability outcomes. Exclusive neighbourhoods had more canopy and better
preventive care \emph{and} more air pollution --- an awkward result the paper does not
resolve. Wholly ecological and cross-sectional. For monitor siting the message is that a
single historical redlining layer is an inadequate stratifier.}
{10.1073/pnas.2422760122}

\paperentry{Lai et al.\ 2026}
{Monte Carlo model for wind-driven PM emissions from industrial granular materials}
{Atmosphere}
{First quantitative evaluation of a saltation-and-sandblasting emission model for wind
erosion of granular stockpiles, with the Monte Carlo impact count fixed rather than run
to convergence. Applied to lead and zinc sulphide concentrates at a Sardinian plant,
cutting simulated impacts from \textbf{10{,}000 to 500} changed the median emission by a
small relative difference at a twentyfold saving, and the revised model reproduced
emission magnitude, sandblasting efficiency and their friction-velocity dependence.
Validation is internal and physical-consistency based --- no measured emission flux is
compared against, which is the central weakness. Fugitive stockpile emission is exactly
the source class regulatory networks miss and fenceline arrays are asked to capture.}
{10.3390/atmos17090891}

\paperentry{Groenen et al.\ 2026}
{Characteristics and trend of wildfire activity in the Southern Hemisphere}
{Sci Total Environ}
{Two decades of MODIS active-fire detections across Southern Africa, South America and
Oceania. Seasonality is strongly heterogeneous: Southern Africa shows a stable 20-week
May--October season with low inter-annual variability, South America a 14-week
July--October season with high variability, and \textbf{Oceania no clear season at all},
which the authors read as co-existing fire regimes. Active-fire counts are a detection
proxy, not emitted mass --- insensitive to smouldering, biased by overpass timing and
cloud --- so no PM emission inventory follows directly. Note that the abstract's stated
period (2004--2024) disagrees with the title's (2003--2022); check the paper before citing
either.}
{10.1016/j.scitotenv.2026.182240}

\paperentry{Shukla et al.\ 2026}
{Organic vapors from savannah and boreal fire emissions under photochemical and dark aging}
{Atmos Chem Phys}
{Controlled burns of understudied African savannah and European boreal fuels sampled by
high-resolution PTR-MS, then aged photochemically and in the dark in a 29\,m$^3$ Teflon
chamber. Total primary organic-vapour emission factors spanned \textbf{69--161\,g/kg} by
fuel type. Photochemical ageing depleted furanics, phenolics and oxygenated aromatics and
enriched carbonyls and small O-containing species; dark low-NO$_x$ ageing produced almost
no compositional change, and clustering attributed \textbf{73\% of group-level
compositional variance to oxidation regime}. Chamber burns of prepared fuels do not
reproduce field moisture, combustion phase or dilution, so these are upper-bound
laboratory factors. Any smoke PM$_{2.5}$ forecast assuming a fixed downwind secondary yield
is contradicted here.}
{10.5194/acp-26-12977-2026}

\band{Occupational \& indoor exposure\hfill 3 records}{Clay}

\paperentry{Bellacov et al.\ 2026}
{Diesel-related pollutant migration in fire stations measured with low-cost direct-reading
monitors}
{J Occup Environ Hyg}
{Seven consecutive days of minute-level monitoring at fixed engine-bay and adjacent
occupied locations in three fire stations, aggregated to daily means and fitted by
hierarchical mixed-effects models. Coarse PM, NO$_2$, VOCs and CO$_2$ were significantly
associated with engine activity, source distance and station identity, but the
\emph{direction} varied by pollutant --- and \textbf{PM$_{2.5}$ and PM$_1$ were not
significantly associated with any of them}. Measurable concentrations beyond the bays
indicate migration or shared indoor sources; this design cannot separate them. Daily
aggregation of minute data discards precisely the transient peaks that dominate diesel
exposure, so the null on the fine fractions reads as a resolution artefact rather than an
absence.}
{10.1080/15459624.2026.2722264}

\paperentry{Fermoso et al.\ 2026}
{Integrated hospital IAQ monitoring: continuous sensing with targeted chemical
characterisation}
{Atmosphere}
{Thirty months of continuous low-cost sensor monitoring across hospital microenvironments
(June 2023 -- December 2025) yielding \textbf{1.79 million raw records, 1.42 million
retained indoor measurements and 264{,}311 hourly aggregates}, paired with targeted VOC
and aldehyde sampling. Distinct recurrent signatures emerged per area: CO$_2$ tracked
occupancy and ventilation demand, while formaldehyde response was persistent and seasonal
in a pattern consistent with indoor sources plus ventilation heterogeneity. The
formaldehyde inference rests on a low-cost channel with known cross-sensitivity, and the
campaign sampling is too sparse to validate it continuously. The \textbf{79\% retained
fraction} is the number to take away when sizing a long-horizon deployment.}
{10.3390/atmos17090895}

\paperentry{Yang et al.\ 2026}
{Organic contaminants of emerging concern in indoor air: global distribution and hazard
burden}
{Environ Sci Technol}
{Seven OCEC classes spanning over 130 compounds synthesised across residences, offices,
schools and commercial spaces. Microenvironments separate cleanly --- phthalate
alternatives dominate residential air to $2.08\times10^4$\,ng/m$^3$, PFAS schools to
230\,ng/m$^3$, organophosphate esters offices to 159\,ng/m$^3$ --- and the US, Germany,
Sweden and China rank highest on inhalation hazard, which the authors say plainly may
reflect monitoring intensity rather than contamination. \textbf{Microplastics were
excluded from the hazard index because units and toxicity metrics are irreconcilable}:
a candid admission that the particulate fraction of indoor emerging contaminants cannot
currently be placed on a common scale with the rest.}
{10.1021/acs.est.6c06332}

\band{Burden, policy \& mitigation\hfill 2 records}{Coral}

\paperentry{Jiang et al.\ 2026}
{Ambient particulate matter pollution and stroke burden}
{Med Gas Res}
{GBD 2021 stroke deaths and DALYs attributable to ambient particulate matter, decomposed
by region, sex and SDI over 1990--2021 with EAPC trends and concentration and slope indices
of inequality. Absolute deaths rose \textbf{107\%} and DALYs \textbf{101\%} while
age-standardised rates fell --- the familiar decomposition of ageing and population growth
against declining risk. Middle-SDI locations carry \textbf{46\% of the global burden} and
exceed high-SDI on both age-standardised measures. The concentration index suggests
improving relative inequality while the slope index worsened against 1990, and the paper
does not reconcile the two. Attributable burden is a modelled quantity resting on the
exposure surface and the chosen CRF; this is not independent epidemiological evidence.}
{10.4103/mgr.MEDGASRES-D-25-00188}

\paperentry{Liu et al.\ 2026}
{Spatiotemporal evolution of PM$_{2.5}$ and its health and economic effects in Central
China}
{Environ Monit Assess}
{Henan, Hubei and Hunan analysed with global Moran's I, centroid migration and BenMAP,
with Henan carried forward for valuation over 2017--2023. Spatial clustering is strong and
stable (\textbf{Moran's I 0.524--0.738, all $p<0.001$}), high-value clusters sit in Henan
and around Wuhan and Changsha, and the concentration centroid moved only
\textbf{2.86--24.37\,km} per year. The authors explicitly \emph{reject} a persistent
southwestward migration because most displacements fall below their noise threshold ---
a welcome negative result in a literature prone to over-reading centroid drift.
Attributable premature deaths declined with a late rebound, but a rising value of a
statistical life meant monetised benefit grew more slowly than health benefit. BenMAP
output inherits its CRF and baseline-rate assumptions entirely.}
{10.1007/s10661-026-15911-z}

\clearpage

\band{Respiratory \& allergic\hfill 1 record}{Sage}

\paperentry{Kim et al.\ 2026}
{Heat exposure and COPD morbidity and mortality: systematic review and meta-analysis}
{Int J Environ Health Res}
{Thirty-six observational studies pooled in a three-level random-effects model that
handles non-independent estimates. Heat raised overall COPD risk
\textbf{(RR 1.15, 95\% CI 1.09--1.21)}, with mortality \textbf{(1.26, 1.18--1.34)} well
above morbidity \textbf{(1.05, 1.01--1.09)}. Heat-wave and threshold definitions both gave
RR 1.21 against 1.05 for continuous temperature metrics, and the continuous metrics were
also the least heterogeneous --- so definitional choice, not biology, drives much of the
spread. The finding that earns this record its place here is negative:
\textbf{adjusting for air pollution did not substantially alter the associations}, which
argues against PM confounding of heat--COPD associations while saying nothing about
interaction. No publication-bias or exposure-misclassification assessment is reported in
the abstract.}
{10.1080/09603123.2026.2734444}

\band{Mechanistic toxicology\hfill 1 record}{Violet}

\paperentry{Vahdatbin et al.\ 2026}
{Is recycling sustainable: the promise and risks of scrap tire valorization}
{iScience}
{Roughly \textbf{1.5 billion end-of-life tires} accumulate annually, and this review argues
that devulcanization, pyrolysis and mechanical grinding each release hazardous species
including VOCs, PAHs, heavy metals, microplastics and 6PPD-quinone. It is qualitative:
no emission factors, exposure concentrations or dose--response relationships are given, so
the risk claim is directional and cannot rank the three routes. Its value here is the
explicit call for high-fidelity monitoring of recycling facilities --- a fenceline
low-cost-network application with essentially no published deployment literature. Tire
wear is already a leading non-exhaust PM source; tire \emph{reprocessing} as a second,
point-source pathway is comparatively unmeasured.}
{10.1016/j.isci.2026.117462}

\clearpage

% ==================================================================== TABLE
\section*{\sffamily\bfseries\color{Deep}Corpus register}
\vspace{-1.5mm}{\color{Amber}\rule{\linewidth}{1.1pt}}\vspace{2.5mm}

{\fontsize{6.7}{8.4}\selectfont
\setlength{\tabcolsep}{2.4pt}
\begin{longtable}{@{}>{\RaggedRight}p{25mm}>{\RaggedRight}p{41mm}>{\RaggedRight}p{27mm}%
>{\RaggedRight}p{23mm}>{\RaggedRight}p{18mm}>{\RaggedRight}p{15mm}@{}}
\toprule
\textsf{\textbf{First author}} & \textsf{\textbf{Focus}} & \textsf{\textbf{Journal}} &
\textsf{\textbf{Design}} & \textsf{\textbf{Metric}} & \textsf{\textbf{Region}}\\
\midrule
\endfirsthead
\toprule
\textsf{\textbf{First author}} & \textsf{\textbf{Focus}} & \textsf{\textbf{Journal}} &
\textsf{\textbf{Design}} & \textsf{\textbf{Metric}} & \textsf{\textbf{Region}}\\
\midrule
\endhead
\multicolumn{6}{@{}l@{}}{\cellcolor{Mist}\textsf{\textbf{\textcolor{Deep}{Neuro / mental health}}}}\\[0.6mm]
\href{https://doi.org/10.1001/jamanetworkopen.2026.25015}{Dales} & Wildfire-sourced PM$_{2.5}$ and ED visits for migraine and other primary headache & \textit{JAMA Netw Open} & Case-crossover & PM$_{2.5}$ (wildfire) & Canada\\
\href{https://doi.org/10.1093/ageing/afag215}{Glover} & Umbrella meta-review of urban environmental exposures and cognitive health & \textit{Age Ageing} & Umbrella review / meta-analysis & PM$_{2.5}$ & Global\\
\href{https://doi.org/10.1016/j.nsa.2026.107019}{Goossens} & In utero air pollution exposure and pre-/postnatal brain development & \textit{Neurosci Appl} & Narrative review & PM$_{2.5}$ / PM$_{10}$ / PM$_{\mathrm{coarse}}$ & Global\\
\href{https://doi.org/10.1016/j.isci.2026.116617}{Ma} & PM$_{2.5}$-induced neurotoxicity and depression: exposure risks and mechanisms & \textit{iScience} & Narrative review & PM$_{2.5}$ + components & Global\\
\href{https://doi.org/10.3760/cma.j.cn112150-20250804-00752}{Pan} & Maternal PM exposure and offspring neurodevelopment: effects and mechanisms & \textit{Zhonghua Yu Fang Yi Xue Za Zhi} & Narrative review & PM (bulk) & China\\
\addlinespace[0.9mm]
\multicolumn{6}{@{}l@{}}{\cellcolor{Mist}\textsf{\textbf{\textcolor{Teal}{Cardiovascular \& metabolic}}}}\\[0.6mm]
\href{https://doi.org/10.3760/cma.j.cn112338-20251120-00843}{Huang} & PM$_{2.5}$ and dyslipidaemia in rural Chinese children; dietary effect modification & \textit{Zhonghua Liu Xing Bing Xue Za Zhi} & Cross-sectional (multistage) & PM$_{2.5}$ & China\\
\href{https://doi.org/10.32394/pe/218966}{Kazmierska} & Air pollution and cardiovascular disease: review of current data & \textit{Przegl Epidemiol} & Narrative review & PM$_{2.5}$ / PM$_{10}$ & Global\\
\href{https://doi.org/10.1007/s00125-026-06814-2}{Xi} & Air pollution, built environment and temperature vs. prediabetes phenotypes & \textit{Diabetologia} & Prospective cohort & PM$_{2.5}$ & Germany\\
\href{https://doi.org/10.1161/JAHA.125.046917}{Yang} & Short-term PM$_{2.5}$ and hospitalisation/ED use in atrial fibrillation & \textit{J Am Heart Assoc} & Case-crossover & PM$_{2.5}$ (wildfire-influenced) & USA\\
\addlinespace[0.9mm]
\multicolumn{6}{@{}l@{}}{\cellcolor{Mist}\textsf{\textbf{\textcolor{Sky}{Reproductive \& developmental}}}}\\[0.6mm]
\href{https://doi.org/10.1016/j.jes.2025.10.043}{Chen} & Prenatal air pollution, fetal growth trajectory phenotypes and child neurodevelopment & \textit{J Environ Sci (China)} & Prospective cohort & PM$_{2.5}$ / PM$_{10}$ & China\\
\href{https://doi.org/10.1016/j.ecoenv.2026.120566}{Huang} & Joint prenatal bisphenol + air-pollutant mixture exposure and preschool IQ & \textit{Ecotoxicol Environ Saf} & Prospective cohort (mixtures) & PM$_{2.5}$ / PM$_{10}$ & China\\
\href{https://doi.org/10.1007/s40572-026-00552-8}{Peterson} & Ambient air pollution, PAHs and pubertal development: critical review & \textit{Curr Environ Health Rep} & Critical review & PM$_{2.5}$ components & Global\\
\href{https://doi.org/10.1016/j.envres.2026.125313}{Xie} & Long-term PM$_{2.5}$ and timing of pubertal development in boys and girls & \textit{Environ Res} & Prospective cohort & PM$_{2.5}$ & China\\
\addlinespace[0.9mm]
\multicolumn{6}{@{}l@{}}{\cellcolor{Mist}\textsf{\textbf{\textcolor{Sage}{Respiratory \& allergic}}}}\\[0.6mm]
\href{https://doi.org/10.1007/s00484-026-03289-z}{Bayar Muluk} & Climate change and otolaryngology: airway disease impacts and emission reduction & \textit{Int J Biometeorol} & Narrative review & PM$_{2.5}$ & Global\\
\href{https://doi.org/10.4103/mgr.MEDGASRES-D-25-00164}{Liao} & Short-term air pollution, respiratory symptoms and lung function in asthma & \textit{Med Gas Res} & Panel study & PM$_{2.5}$ / PM$_{10}$ & China\\
\href{https://doi.org/10.3390/epidemiologia7040097}{Sanchez de Prada} & Ambient PM and asthma-coded admissions in children 0-3 y, Spain & \textit{Epidemiologia} & Ecological panel & PM$_{2.5}$ / PM$_{10}$ & Spain\\
\href{https://doi.org/10.1016/j.jhazmat.2026.143076}{Yang} & Allergenic vs. toxic risk of airborne seed hairs compared with ambient PM$_{2.5}$ & \textit{J Hazard Mater} & Field measurement & PM$_{2.5}$ + bioaerosol & China\\
\addlinespace[0.9mm]
\multicolumn{6}{@{}l@{}}{\cellcolor{Mist}\textsf{\textbf{\textcolor{Violet}{Mechanistic toxicology}}}}\\[0.6mm]
\href{https://doi.org/10.1016/j.jes.2025.12.074}{Ding} & PM$_{2.5}$ promotes vascular calcification via miR-27b-5p/HES1 signalling & \textit{J Environ Sci (China)} & In vitro + organoid + murine & PM$_{2.5}$ & China\\
\href{https://doi.org/10.1007/s00360-026-01698-5}{Maraschi} & Osmoregulatory disruption and metal bioaccumulation from settleable atmospheric PM & \textit{J Comp Physiol B} & Controlled exposure (ecotox) & Settleable APM & Brazil\\
\href{https://doi.org/10.1016/j.jes.2025.11.018}{Wu} & Macrophage NLRP3 inflammasome in real-ambient PM-induced hepatic glucose dysregulation & \textit{J Environ Sci (China)} & Real-ambient murine + co-culture & Real-ambient PM & China\\
\addlinespace[0.9mm]
\multicolumn{6}{@{}l@{}}{\cellcolor{Mist}\textsf{\textbf{\textcolor{Clay}{Occupational \& indoor}}}}\\[0.6mm]
\href{https://doi.org/10.3389/fmedt.2026.1796828}{Chen} & Nested suction instrument to cut intraoperative surgical-smoke PM$_{2.5}$ & \textit{Front Med Technol} & Bench / device experiment & PM$_{2.5}$ (surgical smoke) & China\\
\href{https://doi.org/10.1093/annweh/wxag055}{Cole} & Determinants of respirable crystalline silica exposure in Australian tunnelling & \textit{Ann Work Expo Health} & Observational exposure study & Respirable dust / RCS & Australia\\
\href{https://doi.org/10.1016/j.jhazmat.2026.143081}{Ding} & Stage-resolved hair Zn-isotope dynamics as a firefighting combustion-exposure tracer & \textit{J Hazard Mater} & Repeated-measures biomarker & Combustion-derived PM & China\\
\href{https://doi.org/10.3390/jox16040119}{Ramsden} & Xenobiotic hazards in aircraft cabin air & \textit{J Xenobiot} & Review + risk economics & Ultrafine particles & Global\\
\addlinespace[0.9mm]
\multicolumn{6}{@{}l@{}}{\cellcolor{Mist}\textsf{\textbf{\textcolor{Amber}{Sensing, forecasting \& instrumentation}}}}\\[0.6mm]
\href{https://doi.org/10.1016/j.jenvman.2026.130515}{Shi} & Adaptive spatiotemporal graph transformer for multi-site multi-step PM$_{2.5}$ forecasting & \textit{J Environ Manage} & Deep-learning model & PM$_{2.5}$ & China / India\\
\href{https://doi.org/10.3390/bios16070380}{Tsuchiya} & Seasonal and indoor-outdoor characteristics of fluorescent aerosol particles in offices & \textit{Biosensors} & Real-time sensor campaign & Bioaerosol / coarse & Japan\\
\addlinespace[0.9mm]
\multicolumn{6}{@{}l@{}}{\cellcolor{Mist}\textsf{\textbf{\textcolor{Amber}{Exposure assessment \& modelling}}}}\\[0.6mm]
\href{https://doi.org/10.1029/2025GH001621}{Solorzano-Araque} & Multiscale exposure modelling: LOTOS-EUROS + mobility surveys + particle toxicity & \textit{GeoHealth} & CTM + mobility integration & PM$_{2.5}$ / PM$_{10}$ & Colombia\\
\addlinespace[0.9mm]
\multicolumn{6}{@{}l@{}}{\cellcolor{Mist}\textsf{\textbf{\textcolor{Coral}{Burden, policy \& mitigation}}}}\\[0.6mm]
\href{https://doi.org/10.1002/hsr2.72886}{Karamzad} & Ischaemic heart disease burden in MENA 1990-2021; PM attributable fraction 28.2\% & \textit{Health Sci Rep} & GBD secondary analysis & Ambient PM & MENA\\
\href{https://doi.org/10.1371/journal.pone.0347989}{Rees} & HEPA filtration to reduce respiratory infection in care homes & \textit{PLoS One} & RCT process evaluation & Filtered PM & UK\\
\href{https://doi.org/10.1016/j.jes.2026.01.009}{Shan} & Policy-driven emission reductions and post-pandemic rebound, high-resolution inventories & \textit{J Environ Sci (China)} & Emission inventory & PM emissions & China\\
\addlinespace[0.9mm]
\multicolumn{6}{@{}l@{}}{\cellcolor{Mist}\textsf{\textbf{\textcolor{Slate}{Other clinical endpoints}}}}\\[0.6mm]
\href{https://doi.org/10.3390/dj14070440}{Lupita} & PM$_{2.5}$, preventive-care capacity and caries/edentulism burden across the EU & \textit{Dent J} & Ecological panel & PM$_{2.5}$ & EU-27\\
\href{https://doi.org/10.1016/j.jdin.2026.06.008}{Park} & Long-term PM exposure and risk of androgenetic alopecia & \textit{JAAD Int} & Nationwide cohort & PM$_{2.5}$ / PM$_{10}$ & South Korea\\
\href{https://doi.org/10.3171/2026.3.PEDS2615}{Sun} & Environmental exposure and area deprivation vs. paediatric hydrocephalus outcomes & \textit{J Neurosurg Pediatr} & Registry cohort + geospatial & PM$_{2.5}$ (satellite-derived) & USA\\
\addlinespace[0.9mm]
\bottomrule
\end{longtable}}


\vspace{4mm}
\begin{keybox}
{\sffamily\bfseries\footnotesize\color{Deep}Provenance and method}\\[1.2mm]
{\fontsize{7.2}{8.8}\selectfont
Entry window \textbf{2026-09-14 $\to$ 2026-09-14}, contiguous with the 13 September issue;
no gap. Built on 15 September under the standing rule that the current day is not opened
before 22:00 local.
\textbf{Sources.} PubMed \texttt{[Date - Entry]}, run as two separate queries --- health
axis and instrumentation axis --- through both the local harvester and the connector:
8 unique candidates, of which the instrumentation query's 2 were already in the health
set. Registry-wide deposits for this entry date: \textbf{2{,}375} (13 September: 4{,}593).
Europe PMC \texttt{CREATION\_DATE} 3. Crossref windowed by ISSN across 25 tracked
non-PubMed journals: \textbf{54 deposits}, 14 in scope. arXiv relevance query returned
60 records, none inside the window. OpenAlex returned HTTP 429 and contributed nothing,
as on every recent run. Consensus sweep on low-cost sensor calibration returned 10,
of which 4 were net-new by title and 2 of those were then removed by the DOI gate as
previously shipped under rewritten titles.
\textbf{Screening.} 16 shipped, 11 rejected (9 out of scope, 2 duplicates), \textbf{6 held}
--- DOIs verify but no abstract exists in Crossref, Europe PMC, Scholar Gateway or
Consensus, and these are kept out of \texttt{seen.json} so a later issue can carry them.
\textbf{Trial watch.} One registry record updated in window (NCT07817836, University of
Pavia, urban air quality and work stress vs.\ inflammatory cytokines and DNA methylation,
$n=151$, observational, completed); accumulated to \texttt{trials.json} and surfaced in the
weekly rollup, not here.
\textbf{Labels.} Subtopic, design, particle-metric and geography are assigned from title,
abstract, keywords and MeSH; assignment is single-label by dominant contribution and is a
judgement, not a controlled vocabulary. Sixteen new design labels were added to the
classifier in the same edit that wrote them onto records.
Effect estimates are transcribed verbatim from abstracts. All DOI links resolve to the
publisher of record except \texttt{10.4103/mgr.MEDGASRES-D-25-00188}, which PubMed
reports but which is not yet deposited with Crossref or Handle.
Content derived from PubMed and Crossref metadata; abstracts remain the copyright of
their respective publishers.\par}
\end{keybox}

\end{document}
